\noindent{\LARGE Circles}\hfill{\large Lesson 1}\par\vspace{3pt}
\noindent\rule{\linewidth}{0.6pt}\par\medskip

\colorbox{panel}{%
  \parbox{\dimexpr\linewidth-2\fboxsep\relax}{%
    \small\textbf{Teacher's copy: overview.}
    \begin{itemize}[leftmargin=1.6em,itemsep=1pt,topsep=2pt]
      \item \textbf{Aims.} Students use the terms of a circle correctly. They explain why equal chords subtend equal angles at the centre, and the converse. They explain why the perpendicular from the centre bisects a chord, and the converse. They find chord lengths and distances with the Pythagoras theorem.
      \item \textbf{Syllabus.} CBSE Class IX 2026--27, Unit IV Geometry (25 of 80 marks). The syllabus asks students to \emph{explain why} each result holds. The explanations here use congruence (SSS, SAS, RHS).
      \item \textbf{Method.} Students discover each result by solving a short problem first. Give 3--4 minutes, then discuss the answers and state the theorem. Students need a ruler and a compass.
      \item \textbf{Timing.} Section 1: 5 min. Section 2: 15 min. Section 3: 20 min. Practice: 15 min. Exit question: 5 min.
      \item \textbf{Next lesson.} The circle through three points and the circumcircle of a triangle.
    \end{itemize}}}

\section{Parts of a circle \normalfont\small(5 min)}

Start with familiar circles: the wheel of a bullock cart and the Ashoka Chakra on the national flag. Ask: what is the same about every spoke of a wheel?

\begin{definition}
A \textbf{circle} is the set of all points in a plane at a fixed distance from a fixed point. The fixed point is the \textbf{centre}. The fixed distance is the \textbf{radius}.
\end{definition}

\begin{itemize}[leftmargin=1.8em,itemsep=1pt]
  \item A \textbf{chord} is a line segment that joins two points on the circle.
  \item A \textbf{diameter} is a chord through the centre. Diameter $=2\times$ radius.
  \item An \textbf{arc} is a part of the circle. A chord cuts the circle into a \textbf{minor arc} and a \textbf{major arc}.
  \item A \textbf{segment} is the region between a chord and an arc.
  \item A \textbf{sector} is the region between two radii and an arc.
\end{itemize}

\exfig{
  \coordinate (O) at (0,0);
  \fill[black!12] (O) -- (20:2.4) arc (20:75:2.4) -- cycle;
  \fill[black!28] (200:2.4) arc (200:290:2.4) -- cycle;
  \draw (O) circle (2.4);
  \draw (200:2.4) -- (290:2.4);
  \draw (O) -- (20:2.4); \draw (O) -- (75:2.4);
  \draw (150:2.4) -- (330:2.4);
  \fill (O) circle (1.5pt);
  \node[below=2pt] at (O) {$O$};
  \node[left] at (200:2.4) {$P$};
  \node[below] at (290:2.4) {$Q$};
  \node[left] at (150:2.4) {$B$};
  \node[right] at (330:2.4) {$C$};
  \node[above right] at (20:2.4) {$A$};
  \node[above] at (75:2.4) {$D$};
  \node[font=\footnotesize] at (47:1.6) {sector};
  \node[font=\footnotesize,align=center] at (245:2.0) {segment};
  \node[right,font=\footnotesize] at (2.9,0.6) {$OA$: radius};
  \node[right,font=\footnotesize] at (2.9,0.1) {$BC$: diameter};
  \node[right,font=\footnotesize] at (2.9,-0.4) {$PQ$: chord};
}

\facilitatornote{Ask students to draw a circle and name each part on their own drawing. Students often confuse a segment with a sector. A sector always has two radii as its sides; a segment has a chord.}

\section{Equal chords and angles at the centre \normalfont\small(15 min)}

\textbf{Solve first.}
\begin{enumerate}[label=\textbf{P\arabic*.},leftmargin=2.2em,itemsep=3pt]
  \item A circle has centre $O$ and radius $5$ cm. Chords $AB$ and $CD$ are each $6$ cm long. Draw a rough diagram and join $OA$, $OB$, $OC$ and $OD$.
  \begin{enumerate}[label=(\alph*),itemsep=1pt]
    \item Write the three pairs of equal sides in $\tr AOB$ and $\tr COD$.
    \item Are the two triangles congruent? Name the rule.
    \item What can you say about $\angle AOB$ and $\angle COD$?
  \end{enumerate}
  \item In another circle with centre $O$ and radius $5$ cm, $\angle POQ=\angle ROS=60^\circ$.
  \begin{enumerate}[label=(\alph*),itemsep=1pt]
    \item What kind of triangle is $\tr POQ$? Give a reason.
    \item Find the lengths of the chords $PQ$ and $RS$.
    \item What do you notice about the two chords?
  \end{enumerate}
\end{enumerate}

\begin{answerblock}
\ans{P1.} (a) $OA=OC=5$ cm and $OB=OD=5$ cm (radii); $AB=CD=6$ cm (given). (b) Yes, $\tr AOB\cong\tr COD$ by SSS. (c) $\angle AOB=\angle COD$ (corresponding parts of congruent triangles).
\ans{P2.} (a) $OP=OQ=5$ cm, so $\angle OPQ=\angle OQP$. They add up to $180^\circ-60^\circ=120^\circ$, so each is $60^\circ$. The triangle is equilateral. (b) $PQ=5$ cm and, in the same way, $RS=5$ cm. (c) The chords are equal.
\end{answerblock}

\facilitatornote{Let students work on P1 and P2 before you state Theorems 1 and 2. P1 leads to Theorem 1 and P2 leads to its converse.}

\begin{theorem}
Equal chords of a circle subtend equal angles at the centre.
\end{theorem}

\exfig{
  \coordinate (O) at (0,0);
  \coordinate (A) at (100:2.2); \coordinate (B) at (170:2.2);
  \coordinate (C) at (240:2.2); \coordinate (D) at (310:2.2);
  \draw (O) circle (2.2);
  \draw (A) -- (B); \draw (C) -- (D);
  \draw (O) -- (A); \draw (O) -- (B); \draw (O) -- (C); \draw (O) -- (D);
  \pic[draw,angle radius=9pt] {angle=A--O--B};
  \pic[draw,angle radius=9pt] {angle=C--O--D};
  \draw[thin] ($(A)!0.5!(B)!0.1!90:(B)$) -- ($(A)!0.5!(B)!0.1!-90:(B)$);
  \draw[thin] ($(C)!0.5!(D)!0.1!90:(D)$) -- ($(C)!0.5!(D)!0.1!-90:(D)$);
  \fill (O) circle (1.3pt);
  \node[right=2pt] at (O) {$O$};
  \node[above] at (A) {$A$}; \node[left] at (B) {$B$};
  \node[below left] at (C) {$C$}; \node[below right] at (D) {$D$};
}

\textbf{Why it works.} In $\tr AOB$ and $\tr COD$:
$OA=OC$ (radii), $OB=OD$ (radii) and $AB=CD$ (given).
So $\tr AOB\cong\tr COD$ (SSS). Hence $\angle AOB=\angle COD$ (corresponding parts of congruent triangles).

\begin{theorem}[Converse]
If two chords of a circle subtend equal angles at the centre, then the chords are equal.
\end{theorem}

\textbf{Why it works.} In $\tr AOB$ and $\tr COD$: $OA=OC$, $\angle AOB=\angle COD$ and $OB=OD$. So $\tr AOB\cong\tr COD$ (SAS). Hence $AB=CD$.

\begin{example}
In a circle with centre $O$, chords $PQ$ and $RS$ are equal and $\angle POQ=80^\circ$. Find $\angle ROS$ and $\angle ORS$.
\end{example}

\exfig{
  \coordinate (O) at (0,0);
  \coordinate (P) at (60:2); \coordinate (Q) at (140:2);
  \coordinate (R) at (220:2); \coordinate (S) at (300:2);
  \draw (O) circle (2);
  \draw (P) -- (Q); \draw (R) -- (S);
  \draw (O) -- (P); \draw (O) -- (Q); \draw (O) -- (R); \draw (O) -- (S);
  \pic[draw,angle radius=9pt,"$80^\circ$",angle eccentricity=2] {angle=P--O--Q};
  \fill (O) circle (1.3pt);
  \node[right=3pt] at (O) {$O$};
  \node[above right] at (P) {$P$}; \node[above left] at (Q) {$Q$};
  \node[below left] at (R) {$R$}; \node[below right] at (S) {$S$};
}

\textbf{Given:} $PQ=RS$ and $\angle POQ=80^\circ$.\\
\textbf{To find:} $\angle ROS$ and $\angle ORS$.
\begin{enumerate}[leftmargin=1.8em,itemsep=1pt]
  \item $\angle ROS=\angle POQ=80^\circ$ (equal chords subtend equal angles at the centre, Theorem 1).
  \item $OR=OS$ (radii), so $\tr ORS$ is isosceles and $\angle ORS=\angle OSR$.
  \item $\angle ORS+\angle OSR=180^\circ-80^\circ=100^\circ$ (angle sum of a triangle), so $\angle ORS=50^\circ$.
\end{enumerate}
\textbf{Answer:} $\angle ROS=80^\circ$ and $\angle ORS=50^\circ$.

\facilitatornote{Stress that the angles are at the \emph{centre}. Ask: ``Would the result hold if the angles were measured at a point on the circle?'' Leave the question open for a later lesson.}
